\subsection{Finite contracts for the constructed families}\label{ta:constructed_contracts}
These definitions specify the inputs behind Table~\ref{tab:rq2-family-counts} and the Rolling grid. They describe finite models, not implementations of deployment platforms. Component versions have distinct tags even when their transition graphs coincide. Unlisted plant and endpoint-controller edges are absent; unlisted monitor labels stutter, and errors are absorbing. Ordinary events are controllable unless explicitly called UC. There is no explicit command precedence. OLD/NEW monitors are empty except in Rolling+Audit; the stated interval monitors apply until handover. The fixed endpoint graphs determine every reachable old entry and the listed loadable new targets. The main paper defines quiescence, command execution and the compared merges.

\paragraph{Rolling$(n,m)$ and block size $k$.}
For $2\leq n\leq6$, $1\leq m<n$, each old component has one ready state with a $serve_i$ loop. Its new version has booting and ready states, with UC $ready_i$ from booting to ready and $serve_i$ only at ready. A local transfer maps old ready to new booting. The interval monitor starts at $r=n$, decrements $r$ by the transferred block size, increments it on $ready_i$ (capped at $n$), and errors below $m$. Both fixed endpoints are all-ready tuples whose one-state controllers admit all $serve_i$ loops; these supply one entry and one target. Consecutive blocks of at most $k$, including the final remainder, are transferred atomically. The 15 $(n,m)$ pairs compare $k=1$ with $k=n$; the separate 70-cell grid varies $1\leq k\leq n$. Proposition~\ref{prop:rolling-threshold} proves the threshold for these inputs.

\paragraph{Canary$(n,m)$.}
The same 15 $(n,m)$ values use one old healthy state $H$ and new states $H,H_P,B_P,B$. Only $H$ serves. A transfer has two adversarial outcomes, $H_P$ and $B_P$; UC reports take $H_P$ to $H$ and $B_P$ to $B$. A controllable restart takes $B$ to $H$. The interval monitor is a bit vector initially zero: a broken report sets bit $i$, a restart clears it, and more than $n-m$ set bits cause error. Other labels stutter. This is a \emph{reported-broken} budget; pending physical states are not identified with already reported failures. The saved certificate also checks physical health separately. The all-healthy old/new tuples and one-state serve-only endpoint controllers give one entry and one target. Merging transfers synchronizes their commands and retains every Cartesian-product outcome. The healthy-outcomes-only and no-recovery variants restrict transfer images or remove restart, respectively; their unchanged both-WIN and both-LOSS results delimit the recovery mechanism.

\paragraph{Canary certificate scope.}
The existing separate check passed on every state of all 30 returned WIN certificates: the 15 parameter pairs under Lazy and Direct-Full. It counts replicas in $H$ or $H_P$ against the threshold $m$; $B_P$ and $B$ do not count. Thus it checks model-state health, while $H_P$ still cannot serve. It does not check continuous service, every explored state or runtime behavior. The other 30 saved checks concern LOSS mechanisms, not physical-health success. Fine and merged transfers share capacity reasoning with Rolling, but their reported-health contracts are different; no safety-preserving reduction to the Rolling game is claimed.

\paragraph{DB-Rolling$(3,m)$.}
Here $m\in\{2,3\}$. Each old component has primary $P$, secondary $S$, and election-pending $E$ states; the new version adds booting $B$. Serving loops exist at $P,S,E$. A controllable $stepdown_{i,j}$ synchronously changes $P_i$ to $S_i$ and $S_j$ to $E_j$; UC $elect_j$ changes $E_j$ to $P_j$. A transfer maps only $S$ to $B$, and UC $ready_i$ changes $B$ to $S$. The interval budget starts at three, decreases by the transferred block size, increases by one on readiness (capped at three), and errors below $m$; election does not consume this nonbooting budget. The old endpoint is $(P,S,S)$, the new target $(S,P,S)$, both with serve-only controllers; unreachable padding retains the remaining ordinary alphabet. Thus the contract specifies a change of primary. With $m=2$, a fine policy can change the primary and transfer secondary components successively. The merged command requires all-secondary states, which cannot be quiescent while an election remains pending. With $m=3$, even the first fine transfer violates the budget. These models do not assert uninterrupted write service.

\paragraph{Rolling+Audit.}
This is Rolling$(2,1)$ with controllable $log_{old}$ and $log_{new}$ loops in every local plant state. Each version's audit monitor has unlogged, logged and error states. Its own log sets logged; either serve changes logged to unlogged but errors from unlogged. The other log and all remaining labels stutter. The new initializer selects unlogged on all nine tagged physical tuples. The interval state is $(r,o,a)$, initially $(2,1,0)$: transfers/readiness change $r$ as above, old-stop clears $o$, and new-start sets $a$. Safety requires $r\geq1$ and $o\lor a$, enforcing overlap of the requirements. Each fixed all-ready endpoint controller alternates its own log with either serve; both its unlogged and logged states are entries/targets, and unreachable controller padding retains the remaining ordinary alphabet. These specify the memories, initializer and coverage obligation used in the comparison. Merging the two transfers loses; the single old-stop and new-start cannot be combined further, so merging requirement groups alone is only a renaming.

\paragraph{Threads$(n,B)$.}
For $2\leq n\leq6$ and queue capacity $B\in\{1,2\}$, each worker is idle or busy; component 1 additionally stores queue length $0\leq q\leq B$. UC arrival increments $q$ below capacity. At capacity, backpressure disables arrival whereas saturated arrival is a self-loop. Controllable $dispatch_i$ requires $q>0$ and worker $i$ idle, decrements $q$ and makes that worker busy; for $i>1$ it synchronizes with component 1. UC $finish_i$ returns a busy worker to idle. Old/new plant graphs coincide. Each local transfer is identity on idle states and preserves component 1's queue. There are no monitors. One-state endpoint controllers admit all ordinary labels; their reachable products from the empty, all-idle initial tuple give all old entries and all new targets. These two arrival modes yield 20 pairs. Under backpressure, disabling dispatch lets the finite queue fill and busy workers finish, permitting globally quiescent transfer. Saturated arrival can persist forever and prevents even fine transfer under the contract's quiescence rule. The negative control therefore distinguishes a genuine transfer opportunity from an implicit fairness assumption; it is not a reproduction of a multithreaded runtime.


\subsection{Granularity need not preserve feasibility}\label{ta:granularity-reverse}
The comparisons above do not assert that finer commands always retain every merged solution. The following analytical contract gives the reverse separation; it is not an additional measured instance or part of the 55-pair population.

Both components have one state per version and the shared ordinary alphabet $\{x,y,tick\}$, with $x,y$ UC and $tick$ controllable. Every version has a $tick$ self-loop. The only other loops are: $y$ in old component~1, $x$ in new component~1, $x$ in old component~2, and $y$ in new component~2. All other transitions are absent. Shared labels synchronize, so neither $x$ nor $y$ is enabled in the all-old or all-new product; $x$ loops in new/old and $y$ loops in old/new. Each local transfer maps its unique old state to its unique new state. There are no requirements, initializers or precedence edges. The one-state old and new controllers loop on all ordinary labels. Their closed loops each have one reachable state and a $tick$ loop, hence are safe and deadlock-free; these states supply the entry and target.

With separate transfers, either first transfer reaches a mixed state with an enabled UC self-loop. The remaining transfer is then permanently blocked by the quiescence rule; the environment can repeat the loop forever, so every fine policy loses. Merging both transfers into their single Cartesian-product command instead moves directly from the quiescent entry to the quiescent all-new target, with rank one. Thus fine LOSS/merged WIN is possible. Atomic merging can remove a harmful mixed state, whereas Cell and Rolling show that it can remove a necessary intermediate state. The reported obstructions establish when a specified merge loses, not a universal refinement order between granularities.
